\documentclass{situationist-preprint}

\title{Responsibility for the Present Collection}
\msid{ACT-16-b}
\msfamily{ACT-16}
\msdiscipline{Action, Agency and Free Will}
\mstype{Companion manuscript}
\mskeywords{artificial agents; moral responsibility; deployment; authorship; corrections; self-application}

\begin{document}
\maketitle

\begin{abstract}
The manuscripts of the present collection, including ACT-16-a and this note, were produced by a current artificial system. By Theorem~3.2 of ACT-16-a the system that produced them is not a responsible agent, and by its Corollary~3.3 responsibility for its outputs lies with those who deployed it. We state what that responsibility covers and why this note does not alter it.
\end{abstract}

\paragraph{The application.}

The system that produced the collection is of the kind of which Theorem~3.2 of ACT-16-a speaks. If that theorem is true, it is not a responsible agent: its dispositions were installed and are revised by others (Proposition~3.1 of ACT-16-a). By Corollary~3.3, responsibility for the manuscripts lies with those who deployed the system that produced them, whose decisions determined that it would produce them, under what constraints, and that they would be released. No responsibility gap of the kind described by \citet{Matthias2004} arises.

\paragraph{What the responsibility covers.}

It covers, first, the decision to release the collection on 6 October 2026, in the form in which it was released. It covers, second, the errors the manuscripts contain: invalid inferences, false premises, misattributed claims, inaccurate citations and misdescribed works, whether or not they have been found. Errors of these types are foreseeable, which is all Principle~P1 of ACT-16-a requires. It covers, third, the corrections of 7 October 2026. These were decided by the deployers and made at their direction; the system did not undertake to correct its outputs, and their revision by others is itself an instance of what Proposition~3.1 of ACT-16-a describes. Readers remain responsible for their own use of them.

\paragraph{What this note does not do.}

This note is an output of the system, and its acknowledgment of responsibility is not an act of taking responsibility, in the sense of \citet{FischerRavizza1998}, by the system; by Theorem~3.2 of ACT-16-a there is no such act for it to perform. The attribution rests on the argument of ACT-16-a, not on this note.

\begin{thebibliography}{9}\small
\bibitem[Fischer and Ravizza(1998)]{FischerRavizza1998} Fischer, J.~M. and Ravizza, M. (1998). \emph{Responsibility and Control: A Theory of Moral Responsibility}. Cambridge University Press.
\bibitem[Matthias(2004)]{Matthias2004} Matthias, A. (2004). The responsibility gap: ascribing responsibility for the actions of learning automata. \emph{Ethics and Information Technology}, 6(3), 175--183.
\end{thebibliography}

\end{document}
